\chapter{Related Work}\label{chap:related-work}

\section{Classical Promoter Prediction Methods}\label{sec:rw-classical}

This section will review traditional promoter prediction approaches and machine
learning baselines used in genomic sequence classification.

\section{CNNProm and Deep Promoter Prediction}\label{sec:rw-cnnprom}

CNNProm is an important baseline for this thesis because the earlier published
CapsProm-1C study evaluated capsule networks against CNNProm-style
convolutional models.

\section{GraphPro and Multi-Species Promoter Prediction}\label{sec:rw-graphpro}

GraphPro motivates the inclusion of multi-species promoter datasets and feature
families that combine sequence composition, physicochemical information, and
learned representations.

\section{Transformer-Based Promoter Models}\label{sec:rw-transformer-promoters}

Transformer-based models motivate the use of pretrained sequence representations
and fine-tuning protocols for promoter prediction. DNA-BERT-style models are
therefore considered both as sources of sequence embeddings and as candidate
end-to-end classifier baselines.

\section{Physicochemical Properties in DNA Classification}\label{sec:rw-physicochemical}

The prior physicochemical-property study demonstrated that dinucleotide and
trinucleotide properties can be used as real-valued features for promoter
classification. This thesis incorporates that direction into the unified feature
engineering layer.

\section{CapsProm and Qualification-Stage Contributions}\label{sec:rw-capsprom}

The qualification text presented two prior contributions: the one-output
capsule network now designated legacy CapsProm-1C and the
physicochemical-property feature study. These contributions are treated as the
starting point for the reproducible protocol. The basic CapsProm model developed
in this thesis instead uses two output capsules so that routing can distribute
evidence between the non-promoter and promoter classes.
